% Mục 1. Thuật ngữ, khung giải thích (lưu lượng cache coherence, scheduler, hai loại core) và sáu câu hỏi nghiên cứu.
Nhóm cài đặt bằng C++ mutex lock (TAS, TTAS, MCS), sense-reversing barrier và concurrent counter
(lock, CAS, fetch-and-add). Nhóm kiểm tra tính đúng đắn của các cài đặt và đo khả năng mở rộng trên máy
shared-memory, so với cấu trúc của thư viện. Phần 4 của đề (MPI-3 RMA~\cite{mpi41}) không bắt buộc và nhóm không làm.

\begin{itemize}[nosep]
  \item \emph{Cache line.} Đơn vị (128 byte trên hai máy Apple, 64 byte trên \PCname{LAP1184})
    mà phần cứng giữ nhất quán giữa các core. Nhiều core cùng đọc được, nhưng ghi thì một core phải giữ
    line độc quyền và làm mất hiệu lực bản sao ở core khác.
  \item \emph{Tranh chấp.} Nhiều thread cùng muốn vào critical section. \texttt{OUTSIDE\_WORK}
    $= 0$ là tranh chấp tối đa.
  \item \emph{Vòng lặp đo.} Mỗi vòng, thread lấy lock, gọi \texttt{is\_prime(N)} trong critical
    section, nhả lock, rồi gọi \texttt{is\_prime(N)} thêm OW lần ở ngoài (OW viết tắt
    \texttt{OUTSIDE\_WORK}). $p$ là số thread. Mục Phương pháp đo định nghĩa đầy đủ các ký hiệu, độ
    chậm và độ công bằng.
  \item \emph{TAS, TTAS, MCS.} TAS lặp \texttt{exchange} lên cờ lock. TTAS đọc cờ tới khi rảnh rồi mới
    \texttt{exchange}~\cite{anderson1990spinlock}. MCS xếp hàng, mỗi thread chờ trên cờ
    riêng~\cite{mellorcrummey1991mcs}.
  \item \emph{Episode.} Một lần cả $p$ thread cùng vượt barrier~\cite{hensgen1988barrier}.
  \item \emph{FAA, CAS.} Fetch-and-add cộng và trả giá trị cũ trong một bước. Compare-and-swap chỉ ghi
    nếu giá trị chưa đổi, nếu không thì thất bại và phải thử lại.
\end{itemize}

Mục Kết quả giải thích số đo chủ yếu bằng số lần một cache line đi giữa các core trên đường tới
hạn~\cite{herlihy2020art}, cùng với scheduler và hai loại core của máy. Các hình đánh dấu hai mốc trên
trục số thread, là số P-core và số core của máy (trên \PCname{apple-m5} là \PCval{apple-m5}{pcores} và
\PCval{apple-m5}{ncpu}). P-core là core hiệu năng cao. E-core là core tiết kiệm điện, chậm hơn, và hai
loại nằm ở hai cụm có cache L2 riêng.

\begin{description}[leftmargin=0pt, itemsep=0.2em, font=\normalfont]
  \item[RQ1 (đúng đắn).] Các cài đặt có đúng không, và bộ kiểm tra có đủ nhạy để bắt lỗi cài đặt?
  \item[RQ2 (lock: mở rộng).] TAS, TTAS, MCS mở rộng thế nào theo số thread, \texttt{OUTSIDE\_WORK} và
    $N$, so với lock thư viện?
  \item[RQ3 (lock: backoff và công bằng).] Backoff giúp hay hại, và ở chế độ nào? Lock nhanh nhất có
    công bằng không?
  \item[RQ4 (barrier).] Chi phí một episode tăng thế nào theo số thread, theo cách chờ, so với barrier
    thư viện?
  \item[RQ5 (counter).] Thứ tự hiệu năng giữa FAA, CAS (có/không backoff) và counter dùng lock, và vì
    sao?
  \item[RQ6 (giữa các máy).] Các kết luận trên có giữ nguyên trên máy khác không?
\end{description}

\begin{table}[H]\centering\footnotesize
\begin{tabularx}{\linewidth}{@{}l>{\raggedright\arraybackslash}p{4.4cm}>{\raggedright\arraybackslash}X@{}}
\toprule
RQ & Thí nghiệm & Hình, bảng \\
\midrule
RQ1 & stress, TSan, mutation & Bảng~\ref{tab:corr:mutation}, \ref{tab:corr:summary} (phụ lục) \\
RQ2 & lock $\times$ thread $\times$ OW $\times$ $N$ &
  Hình~\ref{fig:apple-m5:lock_slowdown}, \ref{fig:apple-m5:lock_cs_length} (phụ lục);
  Bảng~\ref{tab:apple-m5:locktime}, \ref{tab:handoff} \\
RQ3 & sáu biến thể backoff, thời gian ngủ thật & Hình~\ref{fig:apple-m5:lock_backoff},
  \ref{fig:apple-m5:lock_fairness} \\
RQ4 & episode/giây, bảy barrier & Hình~\ref{fig:apple-m5:barrier_rate};
  Bảng~\ref{tab:apple-m5:barrier} \\
RQ5 & thời gian, số CAS hỏng & Hình~\ref{fig:apple-m5:counter_time},
  \ref{fig:apple-m5:counter_cas_failures}; Bảng~\ref{tab:apple-m5:countertime} \\
RQ6 & cùng quy trình trên \PCname{C4M7Y670NT} và \PCname{LAP1184} & Hình~\ref{fig:cross:lock},
  \ref{fig:cross:barrier}, \ref{fig:cross:counter}; Bảng~\ref{tab:cross:machines} và các bảng ở
  phụ lục \\
\bottomrule
\end{tabularx}
\caption{Câu hỏi nghiên cứu, thí nghiệm và hình, bảng trả lời.}\label{tab:intro:rqmap}
\end{table}

Đóng góp của báo cáo:
\begin{itemize}[nosep]
  \item cài đặt ba nhóm cấu trúc kèm lập luận memory order cho từng thao tác atomic;
  \item bộ kiểm tra (stress, watchdog, ThreadSanitizer, schedule fuzzing) được đánh giá bằng mutation
    testing;
  \item quy trình đo một lệnh cho đợt đo chính, sinh lại mọi hình và bảng của đợt đo đó, và thêm máy chỉ
    cần thêm tên máy vào danh sách của báo cáo;
  \item so sánh số đo với cận dưới lý thuyết.
\end{itemize}

Mục Nền tảng thí nghiệm mô tả máy đo, và mục Cài đặt trình bày thuật toán. Mục Kiểm tra tính đúng đắn
trả lời RQ1. Mục Phương pháp đo mô tả cách đo, và mục Kết quả và phân tích trả lời RQ2--RQ6. Báo cáo kết
thúc bằng mục Các yếu tố đe dọa độ tin cậy, mục Kết luận và phụ lục.
